\documentclass[11pt]{article}

% ====================================================
% PACKAGES
% ====================================================
\usepackage[margin=1in]{geometry}
\usepackage[T1]{fontenc}
\usepackage{lmodern}
\usepackage{microtype}
\usepackage{setspace}
% \usepackage{parskip}
\usepackage{tcolorbox}
\usepackage{amsmath, amssymb, mathtools}
\usepackage{siunitx}

\usepackage{graphicx}
\usepackage{float}
\usepackage{subcaption}

\usepackage{enumitem}
\usepackage{booktabs}

\usepackage{xcolor}
\usepackage{hyperref}

\usepackage{fancyhdr}
\usepackage{lastpage}

\usepackage{listings}
\usepackage{indentfirst}
\usepackage[justification=centering]{caption}
% \usepackage{algorithm}
% \usepackage{algpseudocode}
% \usepackage{algorithmic}

% \onehalfspacing%


\pagestyle{fancy}
\fancyhf{}
% Header Design
\lhead{\Assignment\ |\ \course}
\rhead{\Name\ |\ \today}
\cfoot{Page \thepage\ of~\pageref{LastPage}}

\usepackage{listings}
\usepackage{xcolor}

\lstset{
    language=Matlab,
    basicstyle=\ttfamily\small,
    keywordstyle=\color{blue},
    commentstyle=\color{gray},
    stringstyle=\color{orange},
    numbers=left,
    numberstyle=\tiny,
    stepnumber=1,
    numbersep=5pt,
    frame=single,
    breaklines=true,
    showstringspaces=false
}


\newcommand{\Name}{Arael Anaya}
\newcommand{\Assignment}{1.1 Exploring a Robotic Mission}
\newcommand{\course}{Space Robotics}


\lstset{
    language=Matlab,
    basicstyle=\ttfamily\small,
    frame=single,
    breaklines=true,
    numbers=left,
    numberstyle=\tiny,
    tabsize=2
}

% makes the hyperlinks look good
\hypersetup{
    colorlinks,
    linkcolor={red!50!black},
    citecolor={blue!50!black},
    urlcolor={blue!80!black}
}

% \setlength{\parindent}{1em}
\usepackage{indentfirst}

\begin{document}

\section*{1. Mission Overview}
Mars has always been of great interest to humans due to the high likelihood of it at one point
holding life. It is the most likely candidate for human colonization and, as such, understanding
it is one of the top priorities for our future in space.

The Mars 2020 Perseverance rover was designed to traverse the harsh environments of Mars as it
attempted to discover signs of ancient microbial life~\cite{rover_components}. The rover was equipped with a number of sensors
and tools that would allow it to search for any signs of life on the planet. Beyond the search for biosignatures, 
the mission was also to help us understand Mars' climate and geology. It used an advanced drilling arm 
to collect rock and soil samples for a future Earth return mission~\cite{rover_components}.

Perseverance also builds directly on the Curiosity mission. While equipped with many of the same sensors
and a similar design, it builds on all lessons learned from Curiosity to advance the rover's
capabilities~\cite{rover_components}. 

\section*{2. Why This Mission Matters to Me}
I find the Perseverance mission particularly interesting because it embodies the idea of continuous 
improvement: every mission should make the next design better. The wheel redesign was my favorite example of this.
Curiosity's aluminum wheels showed significant wear on Mars' sharp basalt rock,
so the Perseverance team redesigned the wheel with a thicker aluminum shell and a new tread pattern
to directly address that failure mode~\cite{rover_components}.

Beyond its iterative design philosophy, Perseverance also represents what I find most interesting
in robotics: autonomous decision-making under tight real-world constraints. Control systems is the
part of the robotics stack I am most passionate about. It's actually what my current internship focuses on. Seeing
Perseverance's dual-computer architecture manage competing tasks in a communication-delayed,
resource-constrained environment makes the concepts I study feel real and impactful~\cite{science_robotics}. The rover's perception pipeline is also an area I am working to understand
more deeply, making the study of its sensors relevant to my own learning.


\section*{3. Topics to Consider}

\subsection*{3.1 Autonomy}
A key challenge for automating the planetary rover is knowing where it is. Early in the mission, the rover
would capture a 360\textdegree{} panorama and a human operator would manually figure out its
position from a map of Mars~\cite{mgl_universetoday}. Advances in computer vision and localization now allow
Perseverance to determine its own position autonomously through a system called Mars Global
Localization (MGL)~\cite{mgl_universetoday}. This lets the rover operate 24/7 toward its goals
without waiting for Earth-based position confirmation~\cite{autonav_jpl}. 

\subsection*{3.2 Sensing and Perception}
Perseverance carries one of the most comprehensive sensor suites ever deployed on a planetary rover.
The Mastcam-Z provides high-resolution stereoscopic imaging and serves as the rover's primary
vision system, while the SuperCam instrument combines a laser, spectrometer, and microphone to
analyze rock composition at a distance~\cite{rover_components}. For close-up analysis, SHERLOC uses an ultraviolet
Raman spectrometer to detect organic compounds and biosignatures directly on rock
surfaces~\cite{rover_components}.

Navigation relies on a separate network of cameras (front and rear Hazcams paired
with NavCams) that feed into AutoNav's terrain modeling pipeline to identify obstacles in real
time~\cite{autonav_jpl}. The tight coupling between the science sensor suite and the navigation perception
hardware reflects a fundamental design constraint in robotics: sensing and autonomy
cannot be architected independently of one another.

\subsection*{3.3 Control Systems}
Perseverance is controlled by two processors operating in parallel, improving
fault tolerance and enabling multiple processes to run simultaneously~\cite{rover_components}. AutoNav holds the record
for the greatest distance driven by a Mars rover without human review: 699.9 meters in a single
drive~\cite{autonav_jpl,science_robotics}.

AutoNav functions as a hierarchical planner. Using MGL's localization output to establish its
position, it generates a 3D map of the surrounding terrain, identifies hazards, and computes a
route before executing it~\cite{autonav_jpl}. A local reactive layer runs alongside this global plan, allowing
the rover to respond to obstacles that appear between planning cycles. This two-layer architecture
is a well-established setup in robust robots, and Perseverance's consistent performance validates its use.

\section*{4. Conclusion}
Perseverance goes beyond just being a tool for exploration. It is a demonstration of what integrated robotic
systems can achieve under extreme real-world constraints. Its autonomous localization stack, layered
control architecture, rich sensor arsenal, and durable mechanical design each reinforce the others,
forming a system that is more capable than any single subsystem could produce alone. As these
capabilities continue to advance in future missions, they help lay the groundwork for the human
exploration of Mars that ultimately depends on knowing the planet can be navigated by our machines first.

\newpage

\begin{thebibliography}{9}

\bibitem{autonav_jpl}
NASA Jet Propulsion Laboratory,
``Autonomous Systems Help NASA's Perseverance Do More Science on Mars,''
\textit{NASA JPL News}, September 21, 2023.
Available: \url{https://www.jpl.nasa.gov/news/autonomous-systems-help-nasas-perseverance-do-more-science-on-mars/}

\bibitem{science_robotics}
V. Verma et al.,
``Autonomous robotics is driving Perseverance rover's progress on Mars,''
\textit{Science Robotics}, vol. 8, no. 80, p. eadi3099, July 26, 2023.
doi: \url{https://doi.org/10.1126/scirobotics.adi3099}

\bibitem{rover_components}
NASA Science,
``Perseverance Rover Components,''
\textit{NASA Science Mission Directorate}, July 22, 2025.
Available: \url{https://science.nasa.gov/mission/mars-2020-perseverance/rover-components/}

\bibitem{mgl_universetoday}
Universe Today,
``NASA's Techno-Wizardry Grants The Perseverance Rover Greater Autonomy,''
\textit{Universe Today}, February 20, 2026.
Available: \url{https://www.universetoday.com/articles/nasas-techno-wizardry-grants-the-perseverance-rover-greater-autonomy}

\end{thebibliography}

\end{document}
